%%%%%%%%%%%%
%
% $Autor: Wings $
% $Datum: 2025-10-07 %
% $Pfad: TemplateSensor $
% $Version: 4250 $
% !TeX spellcheck = en_GB/de_DE
% !TeX encoding = utf8
% !TeX root = filename 
% !TeX TXS-program:bibliography = txs:///biber
%
%%%%%%%%%%%%

% Structure
\chapter{Steckbrücken}
\label{ch:Steckbrücken}

\section{Steckbrücken}
Steckbrücken, welche in der Abbildung \ref{fig:Steckbruecke} zu sehen sind, werden verwendet, um elektrische Verbindungen zwischen verschiedenen 
Komponenten herzustellen, insbesondere auf Steckbrettern oder in Schaltungen, 
in denen vorübergehende Verbindungen benötigt werden. Jede Steckbrücke ist 
flexibel und langlebig und bietet eine einfache Möglichkeit, elektrische 
Verbindungen herzustellen und zu ändern, ohne Lötarbeiten durchführen zu müssen.

\begin{figure}[h]
	\begin{center}
		\includegraphics[width=12cm]{PlugInBridge/PlugInBridge}
		\caption{Steckbrücken}
		\label{fig:Steckbruecke}
	\end{center}
\end{figure}

\subsubsection{Technische Daten}
\begin{description}
	\item[Spannung:] max. 30 V
	\item[Strom:] max. 3 A
\end{description}